\chapter*{Preface}
\addcontentsline{toc}{chapter}{Preface}

Finance is entering a period in which its visible forms are changing faster than its underlying economic problems. Payments are increasingly digital, trading is increasingly automated, credit decisions draw on new forms of data, financial intermediation is moving across institutions and platforms, and artificial intelligence is lowering the cost of producing and processing information. At the same time, leverage, liquidity, incentives, contractual obligations, and systemic interdependence remain central to financial outcomes. Technological change therefore makes a stable analytical structure more necessary, not less.

This book is organized around one first principle:

\begin{quote}
\textbf{Finance is the economic system through which societies organize exchanges across time and uncertainty. It accomplishes this by generating information, designing and enforcing contracts, and valuing future contingent cash flows. Financial markets, institutions, and instruments are different organizational forms through which these three fundamental functions are implemented.}
\end{quote}

The proposition is a point of departure, not a division of finance into three compartments. Its deeper foundation lies in economics. Resources must be allocated across time; future states are uncertain; individuals differ in preferences, endowments, information, and tolerance for risk; and transactions take place in environments where information is costly, incentives are imperfect, and enforcement is never automatic. If these frictions disappeared, much of the institutional structure of finance would lose its economic purpose. The financial system acquires form because real economies depart from that frictionless benchmark.

Information economics is especially important in understanding this departure. Finance is concerned with promises whose consequences lie in the future, and the parties to those promises rarely possess the same information. Borrowers know more about their projects than lenders, managers know more about firms than outside investors, and insured parties often know more about their behavior and risk than insurers. Screening, signaling, monitoring, disclosure, reputation, collateral, governance, and intermediation arise because information is costly and asymmetrically distributed. They are not peripheral details of financial practice; they help explain why particular markets and institutions exist at all.

Information problems are inseparable from incentives. Once resources have been transferred, behavior may change. A borrower can alter risk, a manager can pursue private benefits, an intermediary can exploit informational advantages, and an investor can withdraw liquidity when doing so is individually rational but collectively destabilizing. Financial arrangements therefore evolve not simply to transmit funds, but to shape behavior under uncertainty. Some of these arrangements take contractual form; others are embedded in governance, market design, institutional reputation, regulation, accounting, clearing, settlement, and the allocation of control rights. Contract is important, but it is one element within a broader economic architecture of incentives and institutions.

This perspective also clarifies the place of valuation. Valuation is not an independent foundation of finance; it is the decision problem that emerges once future benefits, costs, and risks must be compared in the present. Interest rates connect consumption and investment across time. Risk premia compensate for bearing uncertainty. Market prices aggregate information imperfectly and, in turn, influence financing and investment. Valuation therefore sits within the larger economic problem of allocating scarce resources when time, uncertainty, information, and incentives matter.

The organization of this book follows that logic. Part I begins with the foundations of the financial system: intertemporal exchange, uncertainty, money, payments, and liquidity. Part II develops the economic frictions that give financial arrangements their form, with particular attention to information, incentives, intermediation, and valuation. Part III examines the architecture and operation of financial markets, including trading venues, exchange-traded and over-the-counter arrangements, price discovery, clearing, settlement, and market liquidity. Chapter 11 therefore treats derivatives primarily as markets and market infrastructure. Part IV examines the institutions through which financial functions are organized, beginning with commercial banks and extending to investment banks, asset managers, insurers, and shadow banking and digital platforms. Chapter 16 builds on Chapter 12 by examining how credit intermediation, funding, collateral, liquidity transformation, and risk bearing can move beyond traditional bank balance sheets. Part V focuses on the construction, valuation, and risk-management use of financial instruments, including forwards, futures, swaps, options, credit derivatives, and structured claims. It therefore develops the instruments introduced in the market context of Part III rather than repeating that market-structure analysis. Part VI turns to the financial system as a whole, examining systemic risk, central banking, regulation, and policy. Its final chapter revisits the economic frictions introduced in Chapter 1 and the systemic vulnerabilities examined in Chapter 21, using AI, digital money, tokenization, and decentralized finance to test how the organization of finance may change while its underlying economic problems persist.

The method of exposition follows the same hierarchy. Economic mechanisms come first; institutional forms, models, and techniques follow from them. Mathematics is used when it makes a relationship precise. Cases, balance sheets, diagrams, and data are used when they reveal mechanisms that prose alone would obscure. Chinese financial practice is placed within a broader economic framework so that differences in institutions can be understood without mistaking institutional form for economic principle.

Artificial intelligence belongs naturally in this framework because it changes the cost of acquiring, processing, and using information. It may alter screening, research, trading, monitoring, and organizational boundaries. But lower information-processing costs do not remove scarcity, uncertainty, incentives, liquidity constraints, legal enforceability, or systemic externalities. In some settings, faster and more homogeneous decision-making may even strengthen feedback and common exposures. The relevant question is therefore not whether AI transforms finance---it already does---but which economic frictions it reduces, which it leaves intact, and which new forms of dependence it creates.

Finance ultimately concerns the organization of economic life when decisions made today depend on an uncertain future. Its institutions change because technology, law, and markets change; its central problems persist because time, uncertainty, information, incentives, and scarcity persist. A durable study of finance should therefore begin not with the names of institutions or instruments, but with the economic forces that make them necessary and the conditions under which they succeed or fail.

\section*{How to Use This Book}

The chapters use a common set of learning components at the end of the chapter. Their selection depends on the subject matter and learning objectives of each chapter; not every component is required in every chapter. When included, the components follow a consistent order: Summary, Mind Map, Key Terms, Financial AI Lab, Discussion Questions, Applied Problems, Data Application, Case Analysis, and Further Reading. This approach keeps the organization predictable while allowing each chapter to include only those components that add substantive value to its analysis and teaching.

That is the perspective from which this book proceeds.
